% !TeX root = ../main.tex
%%% ---------------------------------------------------------------------------
%%% 结论
%%%
%%% 写作提示：结论要与\cref{sec:objective} 的每一条**逐条对应**。
%%% 目的列了 4 条，结论就该能回答这 4 条——批阅时这是最快的检查方式。
%%% 结论里的每个数字都要在前面出现过，不要引入新数据。
%%% ---------------------------------------------------------------------------
\section{结论}
\label{sec:conclusion}

\begin{enumerate}
  \item 在 ResNet-18 + \dataset{} 的常用配置下，三种优化器的测试准确率为
        SGD-M $\qty{94.82(7)}{\percent}$ $>$ AdamW $\qty{94.56(6)}{\percent}$
        $>$ Adam $\qty{93.21(7)}{\percent}$；收敛速度则相反，
        Adam 与 AdamW 达到 \qty{90}{\percent} 各需约 18 与 19 轮，
        SGD-M 需约 32 轮。

  \item 上述差异均超出随机波动：SGD-M 与 AdamW 之差为不确定度的约 3 倍，
        Adam 与 AdamW 之差约为 15 倍（\cref{sec:analysis:significance}）。
        因此「自适应方法收敛快但泛化略差」在本实验条件下成立。

  \item 学习率敏感性方面，SGD-M 在 $\lr$ 偏离 10 倍时准确率下降超过
        \qty{5}{\percent}，Adam 类方法不足 \qty{1.5}{\percent}。
        超参数预算紧张时应选 Adam 类，追求上限且能调参时应选 SGD-M。

  \item 结论的适用边界：仅覆盖 ResNet-18 与 \dataset{}，且各优化器未做
        逐一精细调优（\cref{sec:analysis:error}），不能外推到其他结构或
        数据规模。
\end{enumerate}

若要进一步收紧结论，下一步应当扩大重复次数至 $n = 10$ 以缩小不确定度，
并对每种优化器独立做贝叶斯超参数搜索，排除调优不充分带来的系统误差。
